\subsection*{设备调用的参数类型}

设备调用可传递以下参数：

\begin{itemize}[leftmargin=1.6em]
\item 带形状数据。
\item 标量整数。
\item 目标平台支持的标量浮点值。
\end{itemize}

控制整数可以使用 \co@int@、\co@s32@ 或 \co@u32@。可传递的浮点类型取决于目标平台，例如 \co@float@、\co@double@、\co@half@、\co@bfp16@ 或 \co@half8@。

Tileflow 侧不进行这些浮点值的算术运算，而是把它们传给设备函数用于计算。

下面将几类参数传给设备辅助函数：

\begin{lstlisting}[language=croqtile,caption={向设备函数传递参数}]
__device__ void bar(float *p, int m, int n) {}
__device__ void foo(float *, int, unsigned, double, float) {}
__co__ void foobar(f32 [M, 24] input, int N, float padding) {
  parallel p by 1 {
    shared f32 [14, 7] buffer;
    call bar(input, input.span(0), buffer.span(0));
    call foo(buffer, N, 3, 3.14, padding);
  }
}
\end{lstlisting}

带形状数据传到设备函数时，参数是指向元素的指针，不再携带完整形状。例如 \co@bar@ 的 \co@float *p@ 对应 \co@input.data@；需要的维度另以整数参数传递。设备函数的声明应与这些参数类型匹配。

目标平台没有对应的原生类型时，可以使用 \co@croq::f16@、\co@croq::bf16@ 等类型。
